\documentclass[11pt,a4paper]{article}
\usepackage[utf8]{inputenc}
\usepackage[T1]{fontenc}
\usepackage[spanish]{babel}
\usepackage{hyperref}
\usepackage{geometry}
\geometry{margin=2.5cm}
\usepackage{parskip}

\title{079 hacer todo lo posible desde el propio sesgo, con friccion}
\author{Alejandro Valenzuela Rivera\\\small Universidad Autónoma de Nuevo León}
\date{\today}
\begin{document}
\maketitle

\begin{quote}
\textbf{Sesgo declarado.} Este trabajo se sostiene desde el sesgo declarado de su autor, y no finge neutralidad. Lo que afirma se puede tumbar, y ese es el punto: el falsificador está escrito, no escondido.
\end{quote}

\section*{Resumen}
El guantelete no limita: permite todo lo posible desde el sesgo declarado, con una sola condicion, que haya friccion. Si fuera solo una opinion en grande, Infinity War habria acabado en quince minutos. El proceso nunca acaba, y lo que se deja es el piso para que se sepa que se intento cuando todo salga mal por culpa del autor.

\section{Por qué se sostiene}
La capacidad sin testigo no cambia nada: solo cambia el mundo, y no deja a quien preguntar. La friccion no es un freno al poder: es lo que hace que el poder sea contestable. Y el piso no existe para tener razon: existe para que el que venga sepa que se intento, y con que.

\section{De dónde viene}
La autoridad, en la conversacion del 2026-10-26.

\section{El caso}
Infinity War: si Thanos hubiera tenido friccion, dura quince minutos.

\section{Método}
Un acto que hace todo lo posible desde su sesgo, y deja quien lo conteste.

\section{Límites, y qué lo tumba}
Un acto que hace todo lo posible y no deja a nadie que lo contradiga.

\section{Lo que queda abierto}
\textquestiondown{}Si el proceso nunca acaba, que se deja para que se sepa que se intento?

\section*{Cómo se leyó esta tesis}
Esto se escribió en el repositorio man-made-empathy, del que este candado es una pieza: candados/06 del caso/079 hacer todo lo posible desde el propio sesgo, con friccion.md. La máquina que lo compuso lo hizo desde este mismo archivo, y no agregó nada que no estuviera escrito. Si algo de esto es falso, se dice en el mismo sitio, y se corrige con fecha.

\end{document}
